\section{The lattice Gram operator in the theta spanning construction}
\label{sec:sr-theta-gram}

We isolate the quantitative part of the theta spanning argument and
its precise global kernel input.  The strict diagonal dominance below
is a consequence of a lattice count and finite Heisenberg
orthogonality; it is uniform in the field, the level, and the genus.

Let \(V_\mathbb Q\) have dimension \(\rho\) and a nondegenerate
alternating rational form \(B\).  For a full lattice \(L\), put
\(L^\#=\{v:B(v,L)\subset\mathbb Z\}\), and call \(L\) selfdual if
\(L^\#=L\).  For two selfdual lattices set
\[
 m(L,L')=[L:L\cap L'].
\]
The same index is obtained with \(L,L'\) interchanged, since
selfdual lattices have the same covolume.

\begin{srlemma}[Counting selfdual lattices]
\label{lem:sr-selfdual-count}
Fix a selfdual \(L\).  For every \(m\ge1\), there are at most
\(m^{2\rho}\) selfdual lattices \(L'\) with \(m(L,L')=m\).
The index equals one precisely when \(L'=L\).
\end{srlemma}

\begin{proof}
A free abelian group of rank \(\rho\) has at most \(m^\rho\)
subgroups of index \(m\).  For clarity, the Hermite normal form
count \(a_\rho(m)\) satisfies
\[
 a_\rho(m)=\sum_{d\mid m}d^{\rho-1}a_{\rho-1}(m/d).
\]
The formula follows by fixing the last diagonal pivot in triangular
normal form and its residue entries.  Induction and the bound
\(\#\{d:d\mid m\}\le m\) give
\(a_\rho(m)\le m^{\rho-1}\#\{d:d\mid m\}\le m^\rho\).

For \(L'\) as in the statement let \(J=L\cap L'\).
There are at most \(m^\rho\) possibilities for \(J\).
Both \(L\) and \(L'\) lie in \(J^\#\).
Duality reverses inclusion and preserves the index of an inclusion,
as follows by taking a basis and its \(B\)-dual basis.
Therefore
\([J^\#:L]=[L:J]=m\) and
\([J^\#:L']=[L':J]=m\).
For a fixed \(J\), the lattice \(L'\) is among the subgroups of index
\(m\) of the free group \(J^\#\), of which there are at most
\(m^\rho\).  Multiplication proves the bound.
If \(m=1\), then \(L\subset L'\); duality gives \(L'\subset L\).
\end{proof}

Here is the vacuum calculation producing the exponent \(n/2\).
In the finite Schr\"odinger representation associated with
\(J^\#/J\), the two complementary maximal isotropic groups
\(L/J\) and \(L'/J\) each have order \(m=m(L,L')\).
Their pairing is perfect: an element orthogonal to \(L'/J\)
lifts to \((L')^\#=L'\), and its intersection with \(L/J\) is zero.
In the realization by functions on \(L'/J\), one group acts by
translations and the other by its \(m\) distinct characters.
A unit eigenvector for the character group is a delta function;
a unit eigenvector for the translation group has all coordinates
of modulus \(m^{-1/2}\).
Consequently normalized vacua \(v_L,v_{L'}\), with the compatible
semicharacters, satisfy
\[
 |\langle v_L,v_{L'}\rangle|=m^{-1/2}.
\]
If their semicharacters disagree on \(J\), the inner product is
zero instead: acting by an element of \(J\) gives two different
eigenvalues.  Thus the inequality with \(m^{-1/2}\) always holds.
For \(n\) tensor copies it gives \(m^{-n/2}\).
This argument includes the semicharacter phases; omitting them
would fail at the even part of the Heisenberg extension.

In the level construction a test vector is a sum of \(b=|J_{\rm lev}|\)
translated normalized vacua.  Its labels run through one prescribed
finite level orbit, with unit-modulus character coefficients.
A pairing between two different labels vanishes: a common
Heisenberg eigenoperator distinguishes their residue classes.
For matching labels the same translation acts on both vacua and
preserves their inner product.  Therefore the test vectors
\(\Phi_{L,\Lambda}\), for every positive definite auxiliary
unimodular lattice \(\Lambda\), obey
\begin{equation}\label{eq:sr-test-gram}
 \langle\Phi_{L,\Lambda},\Phi_{L,\Lambda}\rangle=b,\qquad
 |\langle\Phi_{L,\Lambda},\Phi_{L',\Lambda}\rangle|
       \le b\,m(L,L')^{-n/2}.
\end{equation}
The equality of label sets and the distinguishing eigenoperator
are necessary compatibility conditions on the finite level
construction.  In the prescribed construction they are obtained
by the same residue representatives in both test vectors; the
coprimality of the two CM level ideals identifies equality of
their character pairings with equality of their residue classes.
This is the finite pairing calculation, rather than a choice of
independent normalizations for each \(L,\Lambda\).

\begin{srproposition}[A uniform diagonal gap]
\label{prop:sr-gram-gap}
Let \(\mathcal L\) be any set of selfdual lattices in \(V_\mathbb Q\),
let \(\mathcal C\) be a finite nonempty set of auxiliary lattice
classes with positive mass weights \(w_\Lambda\), and suppose
\eqref{eq:sr-test-gram} holds for one common \(b>0\).
If \(n\ge4\rho+4\), define
\[
 C(L,L')=\sum_{\Lambda\in\mathcal C}
       w_\Lambda\langle\Phi_{L,\Lambda},\Phi_{L',\Lambda}\rangle,
 \qquad a=b\sum_\Lambda w_\Lambda.
\]
Then \(a>0\), \(C(L,L)=a\), and
\begin{equation}\label{eq:sr-gram-row-gap}
 \sum_{\substack{L'\in\mathcal L\\L'\ne L}}|C(L,L')|
       \le0.65a.
\end{equation}
\end{srproposition}

\begin{proof}
The diagonal assertion follows from \eqref{eq:sr-test-gram}.
For the remaining sum, group the lattices by \(m(L,L')\).
Lemma~\ref{lem:sr-selfdual-count} gives
\[
 \sum_{L'\ne L}m(L,L')^{-n/2}
 \le\sum_{m=2}^\infty m^{2\rho-n/2}
 \le\sum_{m=2}^\infty m^{-2}<0.65.
\]
The last numerical inequality needs no special value of the zeta
function: the sum from \(2\) through \(10\), plus
\(\int_{10}^\infty t^{-2}\,dt\), is less than \(13/20\).
Apply the triangle inequality to the finite positive mass average.
These estimates also prove absolute convergence of every row.
\end{proof}

\begin{srproposition}[The actual inverse operator]
\label{prop:sr-gram-neumann}
Let \(\mathcal B\) be the Banach space of bounded families
\(F(L,x)\) with the fixed holomorphic weight on the common frame
domain.  Equip it with the supremum of the invariant pointwise
norms.  Let \(C\) satisfy
\eqref{eq:sr-gram-row-gap} and act in the lattice variable:
\[
 (\mathcal CF)(L,x)=\sum_{L'}C(L,L')F(L',x).
\]
Then \(\mathcal C\) is invertible and
\begin{equation}\label{eq:sr-gram-inverse}
 F'=\mathcal C^{-1}F
 =a^{-1}\sum_{k=0}^\infty(-R)^kF,\qquad
 R=a^{-1}\mathcal C-1,\qquad
 \|F'\|\le\frac{\|F\|}{0.35a}.
\end{equation}
The inverse preserves holomorphy, weight, and the simultaneous
automorphy law whenever the original test vectors have that law.
\end{srproposition}

\begin{proof}
The absolutely convergent row estimate gives \(\|R\|\le0.65\).
The geometric series converges in operator norm and multiplication
by \(1+R\) telescopes to the identity, giving
\eqref{eq:sr-gram-inverse}.
Normal convergence on compact subsets preserves holomorphy;
the weight relation is linear and passes through each row sum
and the uniform limit.

More explicitly, let a group \(\Delta\) act on the lattices and
frames and let \(\widetilde\mu(\gamma,L)\) be the common unitary multiplier.
The test-vector identity
\(\omega(\gamma)\Phi_{L,\Lambda}
 =\widetilde\mu(\gamma,L)\Phi_{\gamma L,\Lambda}\)
gives
\[
 C(\gamma L,\gamma L')\widetilde\mu(\gamma,L)
        \overline{\widetilde\mu(\gamma,L')}=C(L,L').
\]
If
\(F(\gamma L,\gamma x)=\widetilde\mu(\gamma,L)^{-1}F(L,x)\),
reindexing the absolutely convergent row shows that
\(\mathcal CF\) has the same law.  Thus every partial sum and
its limit in \eqref{eq:sr-gram-inverse} has that law.
\end{proof}

We state the exact analytic input which turns this inverse into
theta spanning.  This identifies the required version of the
doubled Siegel--Weil identity and the frame slab, without replacing
either by a theorem name.
Choose an orthonormal basis \(P_1,\ldots,P_t\) of the nonzero
harmonic highest-weight space with its Fock inner product.
Columnwise unitary action preserves that inner product, so
\[
 \sum_jP_j(uz)\overline{P_j(uw)}
       =\sum_jP_j(z)\overline{P_j(w)}.
\]
Let \(K(x,y)\) be the result of applying the associated paired
polynomial derivatives to the doubled Gaussian kernel.
Let
\(T_{\Lambda,j}(L,x)=\Theta_{\Phi_{L,\Lambda},P_j}(x)\).
For the common finite data the required kernel identity is
\begin{equation}\label{eq:sr-doubled-kernel-input}
 \kappa\sum_{\gamma\in\Delta}
    C(L,\gamma L')\overline{\widetilde\mu(\gamma,L')}K(x,\gamma y)
  =\sum_{\Lambda\in\mathcal C}w_\Lambda
       \sum_{j=1}^t
       T_{\Lambda,j}(L,x)
       \overline{T_{\Lambda,j}(L',y)},\qquad \kappa\ne0.
\end{equation}
Both the order of \(x,\gamma y\) and the conjugate multiplier are
part of the normalization.

Here the slab is the open set of frames
\[
 S_\delta=\{y:\ |\operatorname{Gram}_\tau(y)_{ab}
                  -\operatorname{diag}(1_p,-1_q)_{ab}|<\delta
                       \text{ for all }\tau,a,b\}.
\]
For sufficiently small fixed \(\delta>0\), the additional
archimedean assertions needed are
\begin{equation}\label{eq:sr-slab-reproduction-input}
 \int_{S_\delta}K(x,y)H(y)\,dy=\rho_0H(x),
 \qquad \rho_0>0,
\end{equation}
for every bounded holomorphic section \(H\) of the fixed weight,
and absolute integrability of \(K(x,\cdot)\) on this slab.
The invariant pointwise norm controls the ordinary absolute value
on the slab by one constant for the fixed weight.  In particular
\[
 \sum_{L'}\int_{S_\delta}
       |C(L,L')K(x,y)F'(L',y)|\,dy<\infty
\]
follows from the row bound and that slab bound.

\begin{srproposition}[Finite theta spanning from the normalized kernel]
\label{prop:sr-gram-theta-span}
Suppose the fixed level data satisfy
\eqref{eq:sr-test-gram},
\eqref{eq:sr-doubled-kernel-input} and
\eqref{eq:sr-slab-reproduction-input}, with the indicated absolute
convergence.  Suppose diagonal action of \(\Delta\) on
\(\mathcal L\times S_\delta\) has a measurable fundamental domain,
up to null sets, and preserves product counting and frame
Lebesgue measure.  Then every bounded holomorphic automorphic
family \(F\) of this weight is a finite linear combination, with
coefficients independent of \(L,x\), of the functions
\(T_{\Lambda,j}(L,x)\).
\end{srproposition}

\begin{proof}
Take \(F'=\mathcal C^{-1}F\).
Reproduction and absolute convergence give
\[
 \sum_{L'}\int_{S_\delta}
       C(L,L')K(x,y)F'(L',y)\,dy=\rho_0F(L,x).
\]
Unfold this integral over the measurable fundamental domain for
the diagonal action.
The automorphy law for \(F'\) contributes
\(\overline{\widetilde\mu(\gamma,L')}\), exactly the factor in
\eqref{eq:sr-doubled-kernel-input}.
Using that identity, the result is an integral of a finite sum
\[
 \kappa^{-1}\sum_{\Lambda,j}
   w_\Lambda\,T_{\Lambda,j}(L,x)
       \overline{T_{\Lambda,j}(L',y)}F'(L',y)
\]
over that fundamental domain.  Its value therefore lies in the
finite-dimensional span of the functions
\(T_{\Lambda,j}(L,x)\).

One need not assume in this last step that every displayed
coefficient integral is individually absolutely convergent.
Here is the requisite finite-dimensional argument.
If an integrable function of \(z\), for every parameter \(a\), is
of the form \(\sum_{j=1}^N e_j(a)t_j(z)\), choose a basis
\(e_{j_1},\ldots,e_{j_s}\) of the finite span of the \(e_j\)'s and
rewrite the sum in that basis.  There are \(s\) sample parameters
\(a_1,\ldots,a_s\) for which the evaluation matrix of this basis
is invertible.  Otherwise all evaluation functionals would span a
proper subspace of the dual and some nonzero basis combination
would vanish at every parameter.
The rewritten coefficient functions are consequently fixed
linear combinations of the \(s\) integrable functions at those
sample parameters.  They are integrable, and integration gives
a fixed linear combination of the basis \(e_{j_i}\)'s.
Apply this to \(a=(L,x)\) and the unfolded integrand.
Since \(\rho_0\kappa\ne0\), division proves the result.
\end{proof}

In applying this proposition, the normalized definite
Siegel--Weil formula of
Section~\ref{sec:sw-boundary} is used in the range \(n>4r\).
Its auxiliary Hermitian space is positive definite, so the
unitary theta quotient is compact at infinity and the genus mass
weights are positive.
The comparison must use the same auxiliary lattice class,
finite test function, doubled rational frame, and self-dual
Fourier measure as in \(T_{\Lambda,j}\).
In the Eisenstein expansion, maximal isotropic subspaces that
are graphs of unitary elements contribute the left side of
\eqref{eq:sr-doubled-kernel-input}.
The level support removes the graphs outside \(\Delta\).
The other maximal isotropic subspaces contain a nonzero rational
isotropic direction; their contributions vanish when the specified
finite test vectors have a nonzero-degree Heisenberg relation in
that direction.
Summing paired derivatives over the orthonormal \(P_j\)'s commutes
with the compact auxiliary unitary average by the displayed
Fock invariance, and the definite theta side becomes the finite
right side of \eqref{eq:sr-doubled-kernel-input}.
This explains both the Gram coefficient and the multiplier and
shows why independent local normalizations would be insufficient.

Sections~\ref{sec:sr-level-phase-proof} and~\ref{sec:sr-frame-slab-kernel}
supply the common finite test-vector norms and phases, the outside-level
support, the nonzero-degree isotropic relations, and the explicit
archimedean slab integrability and reproduction.
The multiplier in this section is their actual test-vector multiplier
\(\widetilde\mu\), which includes the fixed level character.
Subsection~\ref{subsec:sr-invariant-heat-jets} supplies the invariant
jet implication from the scalar identity.
The graph-term computation, passage of the isotropic annihilator
through the nongraph terms, the admissibility of the doubled test
family and identification of its Gaussian kernel remain to be
verified in this application.  In particular
the abstract proposition is not asserted to construct those inputs.
In the frozen source their producers are respectively
\srcid{EisensteinSeriesTermUnitaryGraph},
\srcid{ThetaTestFunctionNoCoinvariantsIsotropic},
\srcid{AdmissibleDoubledHermitianThetaFamily},
\srcid{GaussianKernelIntegrableOnThinSlab}, and
\srcid{SetIntegralCompOfStable}.
The finite lattice count, numerical gap, inverse operator, and
finite-dimensional unfolding implication have been proved above.
